% Topic 1 answer template. Fill this in and hand it back.
%
%   1. Put your name, group and university email in the block below.
%   2. For each problem, keep ONE AI use declaration and comment out the other.
%   3. Write your answers where it says so. Everything you write comes out in
%      blue, so a marker can see at a glance what is yours.
%   4. Build it, and hand in the PDF as well as this file:
%
%        latexmk -pdf -shell-escape problem-set-1-template.tex
%
%      The -shell-escape is for minted, which colours the code. Uploading this
%      one file to Overleaf works too, since Overleaf has minted already.
%
% Each problem starts on a new page. If you need more room, keep typing. The
% boxes grow with the text.
%
% Everything is answered in this file. The only separate file is
% transcript.md, which holds your exchange with the assistant for problem 1.5.
% Hand that file in as well. Put it next to this one before you build, since
% the appendix at the end prints it into your PDF. A marker then reads the
% exchange without opening anything else.

\documentclass[11pt]{article}

\usepackage[letterpaper,top=2.5cm,bottom=2cm,left=2.5cm,right=2.5cm]{geometry}
% T2A is here so that a transcript in Russian prints. T1 is still the default,
% so nothing changes for the rest of the document.
\usepackage[T2A,T1]{fontenc}
\usepackage[utf8]{inputenc}
\usepackage{amsmath, amssymb}
\usepackage[dvipsnames]{xcolor}
\usepackage{mdframed}
\usepackage{minted}
\usepackage{multicol}
\usepackage{enumitem}
\usepackage{fancyhdr}
\usepackage{hyperref}
\hypersetup{colorlinks, allcolors=blue}

% A chat interface sometimes puts an emoji in its reply, and pdflatex has no
% glyph for one. If the build stops on such a character, either delete it from
% transcript.md or give it a printable stand-in here. Say in a comment that you
% did, and what the character was:
%
%   \DeclareUnicodeCharacter{1F4A1}{<U+1F4A1>}

% Everything you write is in this colour. The problem statements stay black.
\colorlet{solution}{RoyalBlue}

% Code is ruled down the left edge, as on the problem sheet. Write your own
% code in a minted block too, inside the boxes below.
\definecolor{coderule}{HTML}{40BA21}
\setminted{frame=leftline, framerule=1.5pt, rulecolor=\color{coderule},
  framesep=8pt, xleftmargin=10pt, linenos, fontsize=\small}

\mdfdefinestyle{answerbox}{linecolor=solution, linewidth=1pt,
  roundcorner=3pt, backgroundcolor=solution!4,
  innertopmargin=6pt, innerbottommargin=6pt,
  innerleftmargin=8pt, innerrightmargin=8pt,
  skipabove=4pt, skipbelow=8pt}
\mdfdefinestyle{reasonbox}{linecolor=solution!45, linewidth=0.6pt,
  roundcorner=3pt,
  innertopmargin=6pt, innerbottommargin=6pt,
  innerleftmargin=8pt, innerrightmargin=8pt,
  skipabove=4pt, skipbelow=10pt}

% --- The three things you write in ------------------------------------------
%
%   \answer{...}                       a short answer, in a box
%   \begin{reasoning}[Justification]   longer text: a derivation, a proof,
%     ...                              an explanation. The label is yours to
%   \end{reasoning}                    choose.
%
\newcommand{\answer}[1]{%
  \par\noindent{\bfseries\color{solution}Answer.}%
  \begin{mdframed}[style=answerbox]\color{solution}#1\end{mdframed}}

\newenvironment{reasoning}[1][Justification]%
  {\par\noindent{\bfseries\color{solution}#1.}%
   \begin{mdframed}[style=reasonbox]\color{solution}}%
  {\end{mdframed}}

% --- AI use declaration ------------------------------------------------------
% Every problem carries exactly one of these. Keep the true one, comment out
% the other. A missing declaration is treated as a missing answer.
\newcommand{\aideclaration}[1]{%
  \par\bigskip
  \noindent{\color{solution!70}\rule{\linewidth}{0.4pt}}\par\vspace{2pt}
  \noindent{\color{solution}\small\textbf{AI use declaration.}\ #1}%
  \par\smallskip}
\newcommand{\aideclnone}{%
  \aideclaration{I have not used AI tools when solving this problem.}}
\newcommand{\aideclused}[2]{%
  \aideclaration{I have used AI (#1) in #2 of this problem, in accordance with
  the course rules and the problem statement.}}

\newcommand{\enc}[1]{\mathsf{#1}}

% Problem 1.2 asks you to underline the redex you contract. Use this for it:
%   \redex{(\lambda x.\, x)\;y} \longrightarrow y
\newcommand{\redex}[1]{\underline{#1}}
\newcommand{\noai}{\textcolor{BrickRed}{\textbf{[No AI]}}}
\newcommand{\aiallowed}{\textcolor{RoyalPurple}{\textbf{[AI allowed]}}}

%%%%%%%%%%%%%%%%%%%%%%%%%%%%%%%%%%%%%%%%%%%%%%%%%%%%%%%%%%%%%%%%%%%%%%%%%%%%%%
% PUT YOUR NAME, GROUP AND UNIVERSITY EMAIL HERE
%
% The name below is Valentin Turchin, who designed Refal and invented
% supercompilation. Replace it with your own.
\newcommand{\studentname}{Mariia Chegodaeva}
\newcommand{\studentgroup}{B24-SD-02}
\newcommand{\studentemail}{ma.chegodaeva@innopolis.university}
%%%%%%%%%%%%%%%%%%%%%%%%%%%%%%%%%%%%%%%%%%%%%%%%%%%%%%%%%%%%%%%%%%%%%%%%%%%%%%

\pagestyle{fancy}
\fancyhf{}
\lhead{\small Solution by \textcolor{solution}{\studentname\ (\studentgroup)
  \texttt{\studentemail}}}
\rhead{}
\cfoot{\small Programming Paradigms Fall 2026, Innopolis University}
\rfoot{\small\thepage}

\begin{document}

\begin{center}
  {\Large\bfseries Topic 1. Problem set --- $\lambda$-calculus}\\[0.5em]
  {\large\color{solution}\studentname\ \quad \studentgroup}\\[0.2em]
  {\small\color{solution}\texttt{\studentemail}}
\end{center}

\medskip
\noindent
Problems 1.1 to 1.4 are \noai. Problem 1.5 marks each of its parts separately.
At a
defense you may be asked to redo any derivation from this sheet on paper.

\bigskip

%%%%%%%%%%%%%%%%%%%%%%%%%%%%%%%%%%%%%%%%%%%%%%%%%%%%%%%%%%%%%%%%%%%%%%%%%%%%%%
\section*{Problem 1.1 ($\alpha$-equivalence)\quad\normalsize\noai}

Twelve terms:

\begin{multicols}{3}
\begin{enumerate}\itemsep1pt
  \item $\lambda x.\, \lambda y.\, x\;(y\;x)$
  \item $\lambda a.\, \lambda b.\, a\;(b\;a)$
  \item $\lambda y.\, \lambda x.\, y\;(x\;y)$
  \item $\lambda p.\, \lambda q.\, q\;(p\;q)$
  \item $\lambda x.\, \lambda y.\, y\;(x\;y)$
  \item $\lambda u.\, \lambda v.\, v\;(u\;v)$
  \item $\lambda x.\, \lambda y.\, x\;(x\;y)$
  \item $\lambda m.\, \lambda n.\, m\;(m\;n)$
  \item $\lambda x.\, \lambda y.\, x\;(y\;z)$
  \item $\lambda a.\, \lambda b.\, a\;(b\;z)$
  \item $\lambda a.\, \lambda b.\, a\;(b\;w)$
  \item $\lambda z.\, \lambda y.\, z\;(y\;z)$
\end{enumerate}
\end{multicols}

\begin{enumerate}[label=(\alph*)]
  \item Group them into $\alpha$-equivalence classes.
  \item Three of the terms above mention $z$. Two of those three are
    $\alpha$-equivalent to each other and the third is not. Say which, and what
    makes the difference.
\end{enumerate}

\aideclnone
%\aideclused{Claude Sonnet 5}{item (a)}

\answer{The classes are: \emph{ \{1, 2, 3 and 12  \}, \{4, 5 and 6\}, \{7, 8\}, \{9, 10\}, \{11\}.} }

\begin{reasoning}[Justification for (b)]
  \emph{Terms 9 and 10 are $\alpha$-equivalent but term 12 isn't $\alpha$-equivalent to them because in terms 9 and 10, z is a free variable, while in term 12, z is bound by the outer abstraction $\lambda z$.\\
  $\alpha$-equivalent is obtained by renaming bound variables. Free variables must therefore remain unchanged}
\end{reasoning}

\clearpage
%%%%%%%%%%%%%%%%%%%%%%%%%%%%%%%%%%%%%%%%%%%%%%%%%%%%%%%%%%%%%%%%%%%%%%%%%%%%%%
\section*{Problem 1.2 (Call by name versus call by value)\quad\normalsize\noai}

Reduce the term
\[
  (\lambda x.\, \lambda y.\, y)\;
  \bigl((\lambda f.\, f\;(f\;z))\;(\lambda w.\, w)\bigr)\;
  c
\]
as far as it goes, twice: once under \textbf{call by value} and once under
\textbf{call by name}. Write one $\beta$-reduction per line and underline the
redex you contract. Then say how many steps each took, and what the longer one spent them on.

\aideclnone
% \aideclused{Claude Sonnet 5}{item (a)}

\answer{Call by value takes \emph{5} steps, call by name \emph{2} steps.}

\begin{reasoning}[Derivation, call by value]
  \emph{
  (\lambda x.\lambda y.y)\;
  \redex{((\lambda f.f(f z))(\lambda w.w))}\; c
  \\
  \to
  (\lambda x.\lambda y.y)\;
  ((\lambda w.w)\redex{((\lambda w.w) z)})\; c
  \\
  \to
  (\lambda x.\lambda y.y)\;
    \redex{((\lambda w.w) z})\; c
  \\
  \to
  \redex{((\lambda x.\lambda y.y) z)}\; c
  \\
  \to
  \redex{((\lambda y.y) c)}
  \\
  \to c  
  \\
  \text{So, call by value} = 5 \(\beta\)-reductions
}
\end{reasoning}

\begin{reasoning}[Derivation, call by name]
  \emph{Let A=((\lambda f.f(f z))(\lambda w.w))
  \\
  Then:
  \\
  \redex{((\lambda x.\lambda y.y)A)} c
  \\
  \to
  \redex{((\lambda y.y)c)}
  \\
  \to c
  \\
  \text{So, call by name }= 2 \(\beta\)-reductions}
\end{reasoning}

\begin{reasoning}[What the longer one bought]
  \emph{Call by value spent three extra steps evaluating the first argument, even though that argument is never used because x does not occur in the body $\lambda$y.y. Call by name avoids this unnecessary work}
\end{reasoning}

\clearpage
%%%%%%%%%%%%%%%%%%%%%%%%%%%%%%%%%%%%%%%%%%%%%%%%%%%%%%%%%%%%%%%%%%%%%%%%%%%%%%
\section*{Problem 1.3 (Church pairs)\quad\normalsize\noai}

Recall the encodings from the lab:
\[
  \enc{tru} = \lambda t.\, \lambda f.\, t
  \quad
  \enc{fls} = \lambda t.\, \lambda f.\, f
  \quad
  \enc{pair} = \lambda f.\, \lambda s.\, \lambda b.\, b\;f\;s
\]
\[
  \enc{fst} = \lambda p.\, p\;\enc{tru}
  \qquad
  \enc{snd} = \lambda p.\, p\;\enc{fls}
\]

\begin{enumerate}[label=(\alph*)]
  \item Write a term $\enc{swap}$ for which $\enc{swap}\;(\enc{pair}\;x\;y)$
    reduces to $\enc{pair}\;y\;x$. Here you may use the aliases.
  \item Write the term that extracts $y$ from
    $\enc{pair}\;x\;(\enc{pair}\;y\;z)$.
  \item Verify your answer to (b) by reduction, expanding every alias, so that
    no $\enc{pair}$, $\enc{fst}$, $\enc{snd}$, $\enc{tru}$ or $\enc{fls}$ is
    left.
\end{enumerate}

\aideclnone
% \aideclused{Claude Sonnet 5}{item (a)}

\answer{(a) $\enc{swap} = {}$\emph{\lambda p.pair(\text{snd p})(\text{fst p})}}

\answer{(b) \emph{fst(snd(\text{pair x}(\text{pair y z})))}}

\begin{reasoning}[Derivation for (c)]
  \emph{
  (\lambda p.p(\lambda t.\lambda f.t))((\lambda p.p(\lambda t.\lambda f.f))((\lambda f.\lambda s.\lambda b.b\;f\;s)x((\lambda f.\lambda s.\lambda b.b\;f\;s)\;y\;z)))
  \\
  \to
  (\lambda p.p(\lambda t.\lambda f.f))((\lambda f.\lambda s.\lambda b.b\,f\,s)\;x\;((\lambda f.\lambda s.\lambda b.b\,f\,s)\;y\;z))(\lambda t.\lambda f.t)
  \\
  \to
  (\lambda f.\lambda s.\lambda b.b\,f\,s)\;x  ((\lambda f.\lambda s.\lambda b.b\,f\,s)\;y\;z)(\lambda t.\lambda f.f)(\lambda t.\lambda f.t)
  \\
  \to
  (\lambda s.\lambda b.b\,x\,s)((\lambda f.\lambda s.\lambda b.b\,f\,s)\;y\;z)(\lambda t.\lambda f.f)(\lambda t.\lambda f.t)
  \\
  \to
  (\lambda b.b\,x\,((\lambda f.\lambda s.\lambda b.b\,f\,s)\;y\;z))(\lambda t.\lambda f.f)(\lambda t.\lambda f.t)
  \\
  \to
  (\lambda t.\lambda f.f)\;x((\lambda f.\lambda s.\lambda b.b\,f\,s)\;y\;z)(\lambda t.\lambda f.t)
  \\
  \to
  (\lambda f.f) ((\lambda f.\lambda s.\lambda b.b\,f\,s)\;y\;z) (\lambda t.\lambda f.t)
  \\
  \to
  ((\lambda f.\lambda s.\lambda b.b\,f\,s)\;y\;z)(\lambda t.\lambda f.t)
  \\
  \to
  ((\lambda s.\lambda b.b\,y\,s)\;z)(\lambda t.\lambda f.t)
  \\
  \to
  (\lambda b.b\,y\,z)(\lambda t.\lambda f.t)
  \\
  \to
  (\lambda t.\lambda f.t)\;y\;z
  \\
  \to
  (\lambda f.y)\;z
  \\
  \to 
  y
  }
\end{reasoning}

\clearpage
%%%%%%%%%%%%%%%%%%%%%%%%%%%%%%%%%%%%%%%%%%%%%%%%%%%%%%%%%%%%%%%%%%%%%%%%%%%%%%
\section*{Problem 1.4 (Testing for zero)\quad\normalsize\noai}

A Church numeral applies its first argument to its second $n$ times:
$\enc{c}_n = \lambda s.\, \lambda z.\, s^n\,z$, where $s^0\,z$ is $z$,
$s^1\,z$ is $s\;z$, $s^2\,z$ is $s\;(s\;z)$, and so on. The booleans are
$\enc{tru}$ and $\enc{fls}$ as in problem 1.3.

\begin{enumerate}[label=(\alph*)]
  \item Write a term $\enc{iszro}$ such that $\enc{iszro}\;\enc{c}_0$ reduces
    to $\enc{tru}$, and $\enc{iszro}\;\enc{c}_n$ reduces to $\enc{fls}$ for
    every $n > 0$. You may \textbf{not} use aliases: write the term out.
  \item Verify it on $\enc{c}_0$ and on $\enc{c}_2$ by full $\beta$-reduction,
    expanding every numeral and boolean. If a written-out term repeats, you may
    name it, provided you say what the name stands for.
  \item In one sentence, why does it work for \emph{every} $n > 0$ and not only
    for the one you checked?
\end{enumerate}

\aideclnone
% \aideclused{Claude Sonnet 5}{item (a)}

\answer{(a) $\enc{iszro} = {}$\emph{\lambda n.n(\lambda x. \lambda t.\lambda f.f)(\lambda t.\lambda f.t)}}

\begin{reasoning}[Derivation for (b)]
  \emph{
    (\lambda n.n(\lambda x.\lambda t.\lambda f.f)(\lambda t.\lambda f.t))(\lambda s.\lambda z.z)
    \\
    \to
    (\lambda s.\lambda z.z)(\lambda x.\lambda t.\lambda f.f)(\lambda t.\lambda f.t)
    \\
    \to
    (\lambda z.z)(\lambda t.\lambda f.t)
    \\
    \to
    \lambda t.\lambda f.t
    \\
    Therefore, iszero\;c_0 \to \lambda t.\lambda f.t \text{ which is tru}
    \\
    (\lambda n.n(\lambda x.\lambda t.\lambda f.f)(\lambda t.\lambda f.t))(\lambda s.\lambda z.s(s z))
    \\
    \to
    (\lambda s.\lambda z.s(s z))(\lambda x.\lambda t.\lambda f.f)(\lambda t.\lambda f.t)
    \\
    \to
    (\lambda z.(\lambda x.\lambda t.\lambda f.f)((\lambda x.\lambda t.\lambda f.f)z))(\lambda t.\lambda f.t)
    \\
    \to
    (\lambda x.\lambda t.\lambda f.f)((\lambda x.\lambda t.\lambda f.f)(\lambda t.\lambda f.t))
    \\
    \to
    (\lambda x.\lambda t.\lambda f.f)(\lambda t.\lambda f.f)
    \\
    \to
    \lambda t.\lambda f.f
    \\
    Therefore, iszero\;c_2 \to \lambda t.\lambda f.f \text{ which is fls}}
\end{reasoning}

\begin{reasoning}[Explanation for (c)]
  \emph{For every n>0, $c_n$ applies $\lambda x.\lambda t.\lambda f.f$ at least once, and this function ignores its argument 
  and returns $\lambda t.\lambda f.f$, whereas $c_0$ applies it zero times and therefore returns the initial $\lambda t.\lambda f.t$}
\end{reasoning}

\clearpage
%%%%%%%%%%%%%%%%%%%%%%%%%%%%%%%%%%%%%%%%%%%%%%%%%%%%%%%%%%%%%%%%%%%%%%%%%%%%%%
\section*{Problem 1.5 (Inlining, audited)}

A retry helper. It reports how many attempts fit in the time allowed.

\begin{minted}{python}
timeout = 30                 # seconds we are willing to spend in total

def make_retry(limit):
    def retry(timeout):      # seconds to allow for one attempt
        spent = 0
        attempts = 0
        while spent < limit:
            attempts += 1
            spent += timeout
        return attempts
    return retry

run = make_retry(timeout)
print(run(10))
\end{minted}

\noindent
It prints \texttt{3}: three attempts of ten seconds fit inside thirty.
\emph{Inlining} \texttt{make\_retry} means replacing
\texttt{make\_retry(timeout)} by its body, with the parameter substituted.

\begin{enumerate}[label=(\alph*)]
  \item \aiallowed\ Ask an assistant to do that inlining, and record what it
    produces.
  \item \noai\ Run its answer. If the program no longer prints \texttt{3}, say
    what went wrong, in the language of the substitution rule. If it prints
    \texttt{3}, read what it changed and say how that avoided the trap. Either
    way the evidence is the code it wrote, not its account of what it did.
  \item \noai\ Give a correct inlining.
  \item \noai\ Write the corresponding $\lambda$-term, of the shape
    $(\lambda \mathit{limit}.\, \lambda \mathit{timeout}.\, \ldots)\;
    \mathit{timeout}$, substitute correctly, and name the step the naive
    inlining skipped.
\end{enumerate}

% Here the second line is usually the true one. Keep whichever applies.
% \aideclnone

\aideclused{ChatGPT-5.5}{item (a)}

% Keep whichever of these two applies to you, and delete the other.
\answer{(a) The transcript is submitted as \texttt{transcript.md}, and printed
in the appendix.}
% \begin{reasoning}[My first attempt at inlining, for (a)]
%   \emph{I used no assistant. My own first attempt was:}
% \end{reasoning}

\answer{(b) Its program prints \emph{3}, where the original prints \texttt{3}.
The inlining is \emph{correct}.}

\begin{reasoning}[Justification for (b)]
  \emph{The assistant avoided variable capture by introducing the fresh variable \_limit before the definition of retry. Thus \_limit keeps the value of the outer timeout (30), while the parameter timeout of retry still receives 10. It therefore did not substitute the free timeout underneath the inner timeout binder, where that occurrence would have become captured }
\end{reasoning}

\begin{reasoning}[A correct inlining, for (c)]
  \begin{minted}{python}
timeout = 30              

def retry(retry_timeout):
    spent = 0
    attempts = 0
    while spent < timeout:
        attempts += 1
        spent += retry_timeout
    return attempts

run = retry

print(run(10))
\end{minted}
\end{reasoning}

\begin{reasoning}[Derivation for (d)]
  \emph{(\lambda limit.\lambda timeout.B(limit,timeout))timeout
$$ \equiv_{\alpha} (\lambda limit.\lambda t.\,B(limit,t))\;timeout $$
\\
\to
_{\beta} \lambda t.B(timeout,t)
\\
\begin{tabular}{p{0.95\textwidth}}
The skipped step is $\alpha$-conversion
($\alpha$-renaming) of the inner bound variable before
capture-avoiding substitution \\
\end{tabular}}
\end{reasoning}

\clearpage
%%%%%%%%%%%%%%%%%%%%%%%%%%%%%%%%%%%%%%%%%%%%%%%%%%%%%%%%%%%%%%%%%%%%%%%%%%%%%%
\section*{Appendix. My transcript for problem 1.5}

\noindent
This is \texttt{transcript.md}, printed as it stands. Hand that file in as
well. The appendix is here so that a marker reading the PDF sees the exchange
without opening anything else.

\medskip
\noindent
Copy every prompt and every reply into it, in the order they happened, and
change nothing else. Chat interfaces have no \emph{export transcript} button,
so this is copying and pasting, one message at a time. If the copy loses the
line breaks inside a code block, put them back and say so here.

\bigskip
% The encoding is switched for the listing alone, so that a transcript in
% Russian prints. An English one is unaffected.
%
% If you used no assistant, there is no transcript, and the build says so
% instead of stopping.
\begingroup\fontencoding{T2A}\selectfont

\IfFileExists{transcript.md}%
  {\inputminted[breaklines, fontsize=\scriptsize, baselinestretch=0.9,
     linenos=false]{markdown}{transcript.md}}%
  {\color{solution}\emph{I used no assistant on this problem, so there is no
     transcript. My own first attempt at the inlining is in problem 1.5(a).
  }}
\endgroup

\end{document}
